\section{Threats to Validity}
\label{sec:threats}

\textbf{Language specificity.}
Both implementations are in Go.  The relative performance of
multiprecision vs.\ fixed-width digit generation may differ in languages
with different allocation models (e.g., C/Rust with stack allocation, or
JVM with generational GC).  The \texttt{sync.Pool} optimization in the
\burgerdybvig implementation is Go-specific.

\textbf{Single-architecture performance data.}
Performance benchmarks run on a single x86-64 machine.  While
cross-architecture \emph{determinism} is verified on arm64 VMs (both
implementations produce byte-identical output), arm64 \emph{performance}
benchmarks are deferred due to lack of bare-metal hardware.  VM overhead
makes cross-architecture performance comparison meaningless.  Future
work should include bare-metal arm64 benchmarks on Graviton, Apple
Silicon, or Ampere hardware.

\textbf{Single JavaScript engine oracle.}
Oracle vectors are generated from V8 (via Node.js).  While ECMA-262
\texttt{Number::toString} semantics are normative and engine-independent,
our vectors are verified against only one engine.  This is mitigated by
exhaustive boundary coverage: 286{,}362 vectors including all subnormals
up to $2^{12}$, all binary64 powers of two, subnormal/normal boundary
neighborhoods, \rfcjcs Appendix~B examples with neighbors, and
pseudorandom fill.

\textbf{Workload representativeness.}
Benchmark workloads are synthetic, designed to bracket the performance
space between number-heavy and string-dominant extremes.  Real-world JSON
payloads may have different number density, nesting depth, and key
length distributions.  However, the control group design (number-only
and string-only workloads) bounds the expected range.

\textbf{sync.Pool mitigation.}
The \burgerdybvig implementation uses \texttt{sync.Pool} to reuse
big-integer buffers, reducing allocation cost.  Without pooling, the
performance gap would be wider.  This reflects realistic production usage
but may understate the theoretical algorithmic advantage of fixed-width
approaches.
